
% EJEMPLO DE FIGURE
\begin{figure}[htbp]
    \centerline{\includegraphics{fig1.png}}
    \caption{Example of a figure caption.}
    \label{fig}
\end{figure}

% EJEMPLO DE ECUACION
\subsection{Equations}
\begin{equation}
    a+b=\gamma\label{eq}
\end{equation}

% EJEMPLO DE TABLA
\begin{table}[htbp]
\caption{Table Type Styles}
    \begin{center}
        \begin{tabular}{|c|c|c|c|}
            \hline
            \textbf{Table}&\multicolumn{3}{|c|}{\textbf{Table Column Head}} \\
            \cline{2-4} 
            \textbf{Head} & \textbf{\textit{Table column subhead}}& \textbf{\textit{Subhead}}&          \textbf{\textit{Subhead}} \\
            \hline
            copy& More table copy$^{\mathrm{a}}$& &  \\
            \hline
            \multicolumn{4}{l}{$^{\mathrm{a}}$Sample of a Table footnote.}
        \end{tabular}
    \label{tab1}
    \end{center}
\end{table}


\subsection{Figures and Tables}
Place figures and tables at the top and bottom of columns. 
Avoid placing them in the middle of columns. 
Large figures and tables may span across both columns. 
Figure captions should be below the figures; table heads should appear above the tables. 
Insert figures and tables after they are cited in the text. 
Use the abbreviation ``Fig.~\ref{fig}'', even at the beginning of a sentence.

Use 8 point Times New Roman for Figure labels. 
Use words rather than symbols or abbreviations for Figure axis labels to avoid confusing the reader. 
As an example, write the quantity ``Magnetization'', or ``Magnetization, M'', not just ``M''. 
If including units in the label, present them within parentheses. Do not label axes only  with units.


Be sure that the symbols in your equation have been defined before or immediately following the equation. 
Use ``\eqref{eq}'', not ``Eq.~\eqref{eq}'' or ``equation \eqref{eq}'', except at the beginning of a sentence: ``Equation \eqref{eq} is . . .''

Please don't use the \verb|{eqnarray}| equation environment. Use \verb|{align}| or \verb|{IEEEeqnarray}| instead.

{\BibTeX} does not work by magic. It doesn't get the bibliographicdata from thin air but from .bib files. 
If you use {\BibTeX} to produce a bibliography you must send the .bib files. 

\subsection{Some Common Mistakes}\label{SCM}
The word ``data'' is plural, not singular.

The subscript for the permeability of vacuum $\mu_{0}$, and other common scientific constants, is zero not a lowercase letter o.

Commas, semicolons, periods, question and exclamation marks are located within quotation marks only when a complete thought or name is cited, such as a title or full quotation. 
When quotation marks are used, instead of a bold or italic typeface, to highlight a word or phrase, punctuation should appear outside of the quotation marks. 
A parenthetical phrase or statement at the end of a sentence is punctuated outside of the closing parenthesis (like this). (A parenthetical sentence is punctuated within the parentheses.)

A graph within a graph is an ``inset'', not an ``insert''. The word alternatively is preferred to the word ``alternately'' (unless you really mean something that alternates).

Do not use the word ``essentially'' to mean ``approximately'' or ``effectively''.

In your paper title, if the words ``that uses'' can accurately replace the word ``using'', capitalize the ``u''; if not, keep using lower-cased.

Be aware of the different meanings of the homophones ``affect'' and ``effect'', ``complement'' and ``compliment'', ``discreet'' and ``discrete'', ``principal'' and ``principle''.

Do not confuse ``imply'' and ``infer''.

The prefix ``non'' is not a word; it should be joined to the word it modifies, usually without a hyphen.

There is no period after the ``et'' in the Latin abbreviation ``et al.''.

The abbreviation ``i.e.'' means ``that is'', and the abbreviation ``e.g.'' means ``for example''.